% =====================================================================
% NEW CITATIONS NEEDED
%
% -- Essential --
% Tsimpoukelli2021_frozen   Multimodal Few-Shot Learning with Frozen LMs (NeurIPS 2021)
% Alayrac2022_flamingo      Flamingo (NeurIPS 2022)
% Li2023_blip2              BLIP-2 (ICML 2023)
% Liu2023_llava             Visual Instruction Tuning (NeurIPS 2023)
% Hu2022_lora               LoRA (ICLR 2022)
% Nalisnick2019_dgm         Do Deep Generative Models Know What They Don't Know? (ICLR 2019)
% MerrillSabharwal2023      The Parallelism Tradeoff: Limitations of Log-Precision
%                           Transformers (TACL 2023)
% Dziri2023_faithfate       Faith and Fate: Limits of Transformers on Compositionality
%                           (NeurIPS 2023)
% Gekhman2024               Does Fine-Tuning LLMs on New Knowledge Encourage
%                           Hallucinations? (EMNLP 2024, pp. 7765-7784)
% Golkar2023_xval           xVal: A Continuous Numerical Tokenization for Scientific LMs
%
% -- Recommended --
% Zaheer2017_deepsets       Deep Sets (NeurIPS 2017)
% Kang2024_unfamiliar       Unfamiliar Finetuning Examples Control How LMs Hallucinate
% Schick2023_toolformer     Toolformer (NeurIPS 2023)
% Lochner2021_astronomaly   Astronomaly: personalised anomaly detection in astronomy
%
% -- Cut first if tight --
% MerrillSmith2022_saturated  Saturated Transformers are Constant-Depth Threshold
%                             Circuits (TACL 2022)
% Lee2019_settransformer      Set Transformer (ICML 2019)
% Gruver2023_llmtime          LLMs Are Zero-Shot Time Series Forecasters (NeurIPS 2023)
% Gao2023_pal                 PAL: Program-Aided Language Models (ICML 2023)
% Zheng2023_judge             Judging LLM-as-a-Judge with MT-Bench (NeurIPS 2023)
%
% VERIFY BEFORE SUBMITTING: only Gekhman2024, MerrillSabharwal2023, Dziri2023 and
% Golkar2023_xval were checked against a source. The rest are from memory.
% =====================================================================

\section{Introduction}\label{sec:intro}
The integration of Large Language Models (LLMs) into the scientific workflow marks a
paradigm shift in how researchers interface with information. While these models have
demonstrated remarkable proficiency in some scientific tasks, their utility remains
largely confined to the textual and visual domains. A critical frontier in ``AI for
Science'' is the transition from processing scientific text to reasoning directly over
the physical data that defines scientific discovery. Standard LLM
architectures struggle with numerical reasoning \citep{Rahman2025}, in part because
subword tokenization does not preserve the magnitude and continuity of numerical values
\citep{Golkar2023_xval}, creating a gap between natural language and empirical data.

The main motivation of this work is the need to combine physics foundation models and
LLMs in a scientifically useful manner. Each has distinct advantages, and the potential
for knowledge alignment between them has far-reaching outcomes. One is interpretability:
an LLM that operates over both natural language and learned scientific representations
could serve as a `bridge' between humans and machine learning representations, which are
usually hard to interpret. A second is scientific discovery: foundation models trained on
large and diverse datasets encapsulate latent information beyond their downstream tasks,
and that information can support non-trivial population analysis or complex cross-modal
correlations. Unfolding it, however, requires analysing high-dimensional features
directly, whereas the common practice of projecting them to a low-dimensional summary
discards exactly the structure such analyses depend on. Most current uses of LLMs with
scientific data take the form of textual descriptions, or of agents that delegate the
analysis to dedicated models and read off their results. Whether an LLM can perform
scientific reasoning \emph{directly} from learned data representations is therefore still
an open question, and it is the question we study here.

\textbf{Contribution.}
We connect frozen survey encoders to an 8B language model through a trained
projector, adapting the language model itself with either LoRA or full fine-tuning, and evaluate it on two independent astrophysical datasets (Kepler-field stars,
COSMOS-Web galaxies) across per-object and pairwise questions, and on a more challenging 
question for each: the odd member of a set of stars, and the agreement between two
views of a single galaxy.
%
\emph{First}, we find a sharp asymmetry. Reading learned tokens, the model describes
single objects and compares pairs mostly correctly on both datasets, and judges whether a
galaxy's shape and star-formation history agree for $97$--$99\%$ of galaxies. It fails
when the answer requires comparing a whole group: it names the odd star and its property
in a third of groups, falling from $0.51$ with four stars to $0.15$ with $32$. A linear
probe ranks the odd star as well as the catalog values do, so the signal is in the
features and the failure is in using it. Supplying the catalog values as text does not
help: the model either names most of the group or, when it reasons first, does no better
than the model reading the tokens, and a larger general-purpose model (Qwen3.5-27B) does
no better either.
%
\emph{Second}, performing the group computation outside the model and returning it as a
single extra token doubles the stellar score ($0.33 \to 0.66$), keeps it flat across
group sizes, and holds on groups whose odd star was never the odd one in training. This
narrator was also trained differently (LoRA instead of full fine-tuning), and the
matching control has not been run (\Cref{sec:limitations}). On galaxies, where the
narrator already succeeds, the head token adds nothing.
%
\emph{Third}, verified templated prose yields a substantially more factual model than
fluent LLM-written prose on the same task, because the LLM-written text states more than
the input determines.

\section{Related Work}\label{sec:related_work}

\paragraph{LLMs as scientific assistants.}
Advances in LLMs have already had a wide-ranging impact on how we do science. LLM agents
can now review and write scientific papers \citep{Villaescusa-Navarro2025_denario,
Zhuang2025}, act as a `virtual lab' that suggests and implements novel research ideas
\citep{swanson_virtual_2025}, or perform equation discovery \citep{Yang2026ThinkLA}.
However, these agents typically access scientific data indirectly, by writing code to an
external tool (e.g., ESM, AlphaFold-multimer), an approach formalised in the
tool-augmentation literature \citep{Schick2023_toolformer, Gao2023_pal}. Recently there
has been growing interest in predicting physical properties of proteins with LLMs
\citep{Zhuo2024ProtLLMAI,Dotan2026_betadescribe}. That line is
close to ours, but asks a different question: property prediction and textual description
from sequential data, rather than comparative and population analysis over continuous
data.

\paragraph{Scientific foundation models and scientific LLMs.}
Foundation models have become popular across scientific fields and show strong
performance in many domains \citep[e.g.,][]{Leung2024, Parker2025_aion, Smith2024_astropt,
euclid_2025_multimodal, Bardhan2025, Hayes2025_ESM3, bodnar_2025_aurora,
abramson_2024_alphafold3}. Notably, these models do not treat natural language as a
modality, so a clear separation persists between scientific foundation models and LLMs.
The term scientific-LLM usually refers instead to injecting scientific knowledge into
existing LLMs \citep[for reviews, see][]{Hu2025_sci_llm, Eger2025_sci_llm}, most often by
fine-tuning on high-quality text and images to improve reasoning \citep{Xie2023_darwin,
Bai2025_intern_s1, Cherian2024_llmphy}, object description \citep{Mishra-Sharma2024},
search \citep{Koblischke2025}, and chat \citep{Zaman2025}. Such models outperform general
LLMs on domain tasks \citep{haan2025_astrosage}, but remain confined to text and images.

\paragraph{Connecting frozen encoders to frozen language models.}
The architecture we use is standard outside the sciences: a trained projector maps a
frozen encoder's outputs into the embedding space of a frozen language model, which is
adapted with a small number of trainable parameters \citep{Tsimpoukelli2021_frozen,
Alayrac2022_flamingo, Li2023_blip2, Liu2023_llava}. This recipe is
well-established for producing descriptions of a single input. Whether it supports
comparison across several inputs, which is what population-level scientific questions
require, has received little attention.

